\clearpage
\typeout{MIDTERM_BODY_PAGES=\number\numexpr\value{page}-1\relax}
\begin{thebibliography}{6}
\bibitem{jev} TypeSafe, System One API y Choice.
\url{https://docs.typesafe.ai/api/system-one}.
Consulta: 23 de octubre de 2026.
\bibitem{kahn} A. B. Kahn, Topological sorting of large networks,
\emph{Communications of the ACM}, vol. 5, no. 11, pp. 558--562, 1962.
doi: \href{https://doi.org/10.1145/368996.369025}{10.1145/368996.369025}.
\bibitem{mcp} Model Context Protocol, Python SDK, v1.x.
\url{https://github.com/modelcontextprotocol/python-sdk/tree/v1.x}.
Consulta: 23 de octubre de 2026.
\bibitem{mpi} mpi4py, Tutorial e Installation, versión 4.1.2.
\url{https://mpi4py.github.io/mpi4py/stable/html/tutorial.html}.
Consulta: 23 de octubre de 2026.
\bibitem{graham} R. L. Graham, Bounds on multiprocessing timing
anomalies, \emph{SIAM Journal on Applied Mathematics}, vol. 17,
no. 2, pp. 416--429, 1969.
doi: \href{https://doi.org/10.1137/0117039}{10.1137/0117039}.
\bibitem{heft} H. Topcuoglu, S. Hariri y M.-Y. Wu,
Performance-effective and low-complexity task scheduling for
heterogeneous computing, \emph{IEEE Transactions on Parallel and
Distributed Systems}, vol. 13, no. 3, pp. 260--274, 2002.
doi: \href{https://doi.org/10.1109/71.993206}{10.1109/71.993206}.
\end{thebibliography}

\appendices
\section{Reproducción}
Desde la carpeta del proyecto:
\begin{verbatim}
python3 -m venv .venv
source .venv/bin/activate
python -m pip install -r requirements-lock.txt
python -m unittest discover -s tests -v
python scripts/evaluate.py
python scripts/benchmark_mpi.py
\end{verbatim}
\newpage
Para MCP real distribuido:
\begin{verbatim}
mpiexec -n 4 python -m midterm run \
 examples/mcp_plan.json --backend mpi \
 --no-history --output results/mpi_mcp.json
\end{verbatim}
La evaluación local está en \texttt{results/evaluation.json}; las
referencias, en \texttt{experiments/golden.json}. Los tiempos excluyen
análisis Jev y arranque MPI, pero incluyen procesos por tarea y sesiones
MCP cuando corresponde. La plataforma fue macOS arm64, con el mismo
intérprete entre métodos.
\texttt{results/scaling.json} y su CSV identifican doce ensayos MPI
exploratorios. \texttt{results/mpi.json} verifica siete tareas en cinco
ranks; \texttt{results/mpi\_mcp.json}, herramientas MCP en cuatro ranks.
No se midieron energía ni varios nodos.

\section{Formato y entrega}
Las fuentes usan \texttt{IEEEtran}, modo \texttt{conference}, dos
columnas. Las referencias y estos apéndices se separan del cuerpo.
Se conserva \texttt{ID} en el archivo para que los autores añadan
sus identificadores. Q1--Q5 y M1--M4 distinguen implementación
verificada de trabajo planeado.
